\section{Why a one-dimensional $g$ cannot settle this}
\label{sec:inseparable}

The design above pairs each non-bijective map with a bijective one of the same
degree, on the reasoning that a Gaussian process should find two degree-5
polynomials about equally hard, so any difference between them is
attributable to bijectivity. That reasoning is wrong, and the way it fails
rules out the whole approach.

\paragraph{A monotone $g$ preserves the shape of $h$; a folding one does not.}
If $g$ is increasing, $g(h(x))$ has exactly the peaks and valleys that $h(x)$
has, merely relabelled: the objective is a warped copy of the inner function.
If $g$ folds --- turns around at some interior point --- then every time
$h(x)$ sweeps past that point the objective turns around too, and $y(x)$
acquires a feature that $h(x)$ never had.

Counting turning points along a line through the input space:

\begin{center}
\begin{tabular}{llr}
\toprule
Map & & Turns in $y(x)$ \\
\midrule
$h$ itself & --- & 1 \\
\midrule
\texttt{root1}, \texttt{mono3}, \texttt{mono5}, \texttt{exp\_b1}, \texttt{sinh\_b1}
 & all bijective & 1 \\
\texttt{root3}, \texttt{root5} & both non-bijective & 2 \\
\bottomrule
\end{tabular}
\end{center}

\noindent Every bijective map leaves the shape of $h$ untouched; every
non-bijective one adds a fold. The correspondence is exact, and the same holds
along the root-spread family: $s=0$ gives one turn and every $s>0$ gives two or
three.

\paragraph{This is forced, not incidental.} For a scalar $g$,
\[
\text{non-injective}
\iff \text{non-monotone}
\iff g \text{ folds}
\iff y(x) \text{ gains turning points,}
\]
and a stationary GP fits a function with more turning points less well. So in
one dimension, \textbf{destroying information and making the objective harder
to model are the same operation}. They cannot be separated by choosing a
better one-dimensional family, because the implication runs through
monotonicity and admits no exceptions.

The matched pairs therefore do not do what they were built to do, and neither
does the root-spread sweep. Both move injectivity and model difficulty
together.

\paragraph{What the sweep can still establish.} The confound has a direction,
which makes the evidence asymmetric. Spreading the roots raises both the
ambiguity and the oscillation, and the oscillation pushes the gain
\emph{up}. So:

\begin{itemize}\itemsep1pt
\item a gain that \emph{rises} with $\bar n$ is uninterpretable --- either
cause would produce it;
\item a gain that stays \emph{flat} is strong evidence against injectivity
mattering, since it holds against a tailwind.
\end{itemize}

\noindent The degree-2 row is also informative, though not in the way first
intended: with two roots symmetric about the origin, spreading them shifts the
parabola vertically and the affine normalisation removes the shift exactly, so
all five maps in that row are the \emph{same problem} to machine precision
($2\times10^{-16}$). Five independent runs of an identical problem measure the
Monte Carlo noise floor directly, which calibrates how much movement elsewhere
is real.

\paragraph{The escape is more than one intermediate.} With $k\ge2$, $g$ can map
two different \emph{vectors} to the same value without folding along any
direction. The clearest case is $g(\mathbf y)=\sum_i y_i$: the split of the sum
is unrecoverable, yet $g$ is a plane, nothing folds, and $f$ remains exactly a
Gaussian process so the direct model is correctly specified. Measured, the
objective keeps exactly one turning point at every $k$:

\begin{center}
\begin{tabular}{lrrrr}
\toprule
$g$ & $k=2$ & 4 & 8 & 16 \\
\midrule
$\sum_i y_i$ \quad (linear, correctly specified) & 1 & 1 & 1 & 1 \\
$\max_i y_i$ \quad (one kink, no folds) & 2 & 2 & 3 & 6 \\
$\sum_i y_i^2$ \quad (folds) & 1 & 0 & 2 & 3 \\
\bottomrule
\end{tabular}
\end{center}

\noindent Along the top row the number of discarded dimensions grows with $k$
while model difficulty is pinned, which is the separation one dimension cannot
provide. Two caveats keep this from being the final answer. For the linear
family the discarded information is irrelevant to $f$ --- the objective depends
only on the sum --- so the proposition already predicts exactly zero and the
row is a control rather than a discovery, albeit one at the $K$ values of the
real benchmarks. For $\max$ the discarded information \emph{is} relevant, but
its turning-point count still climbs from 2 to 6, so the pinning is partial.

\paragraph{Where this points.} Every two-arm design we have tried runs into the
same wall from a different direction: regret because bijectivity decides where
the optimum can sit, high $\bar n$ because no bijective map can be matched on
difficulty, and now matched degrees because folding and information loss
coincide. The common cause is that Direct differs from Composite in two ways at
once, and no choice of $g$ separates them. A model that holds the prior on
$\mathbf h$ fixed and varies only whether $\mathbf h$ is observed differs in one
way by construction, at any $g$ --- which is the three-arm comparison, arrived
at here from a third independent direction.
